% ==============================================================================
% PREAMBULE — Dossier d'Architecture LockBits
% ==============================================================================
% Ce fichier contient tous les packages, styles, couleurs et configurations
% nécessaires a la compilation du dossier d'architecture.
% ==============================================================================

\documentclass[12pt,a4paper,twoside]{report}

% ==============================================================================
% 1. PACKAGES DE BASE
% ==============================================================================
\usepackage[utf8]{inputenc}               % Encodage UTF-8
\usepackage[T1]{fontenc}                  % Encodage des polices
\usepackage[french]{babel}                % Typographie française
\usepackage{textcomp}                     % Symboles supplémentaires
\usepackage{csquotes}                     % Citations francaises

% ==============================================================================
% 2. MISE EN PAGE
% ==============================================================================
\usepackage[left=2.5cm, right=2.5cm, top=2.5cm, bottom=2.5cm, bindingoffset=0.5cm]{geometry}

% ==============================================================================
% 3. POLICES ET TYPOGRAPHIE
% ==============================================================================
\usepackage{titlesec}                     % Personnalisation des titres

% ==============================================================================
% 4. COULEURS
% ==============================================================================
\usepackage{xcolor}

% Couleurs principales du projet
\definecolor{primaryblue}{RGB}{0, 82, 136}
\definecolor{accentorange}{RGB}{230, 126, 34}

% Couleurs pour les listings de code
\definecolor{codebg}{RGB}{245, 245, 245}
\definecolor{codeframe}{RGB}{200, 200, 200}

% Couleurs semaphoriques
\definecolor{successgreen}{RGB}{39, 174, 96}
\definecolor{warnyellow}{RGB}{241, 196, 15}
\definecolor{errorred}{RGB}{192, 57, 43}
\definecolor{infoblue}{RGB}{52, 152, 219}

% ==============================================================================
% 5. LISTINGS DE CODE
% ==============================================================================
\usepackage{listings}

% Configuration globale des listings
\lstset{
  basicstyle=\ttfamily\small,
  breaklines=true,
  extendedchars=true,
  frame=single,
  backgroundcolor=\color{codebg},
  numbers=left,
  numberstyle=\tiny\color{gray},
  numbersep=5pt,
  keywordstyle=\color{primaryblue}\bfseries,
  commentstyle=\color{green!50!black},
  stringstyle=\color{accentorange},
  showstringspaces=false,
  tabsize=2,
  captionpos=b,
  % Support des caracteres français dans le code
  literate={é}{{\'e}}1 {è}{{\`e}}1 {ê}{{\^e}}1 {ë}{{\"e}}1
           {à}{{\`a}}1 {â}{{\^a}}1 {ù}{{\`u}}1 {û}{{\^u}}1
           {î}{{\^i}}1 {ï}{{\"i}}1 {ô}{{\^o}}1 {ö}{{\"o}}1
           {ç}{{\c{c}}}1 {É}{{\'E}}1 {È}{{\`E}}1 {Ê}{{\^E}}1
           {À}{{\`A}}1 {Â}{{\^A}}1 {Ù}{{\`U}}1 {Û}{{\^U}}1
           {Î}{{\^I}}1 {Ô}{{\^O}}1 {Ç}{{\c{C}}}1
           {€}{{\euro}}1
           {°}{{\textdegree}}1
}

% Definition du langage YAML
\lstdefinelanguage{yaml}{
  keywords={true, false, null, y, n, yes, no, on, off},
  keywordstyle=\color{primaryblue}\bfseries,
  sensitive=false,
  morecomment=[l]{\#},
  morestring=[b]',
  morestring=[b]",
}

% Definition du langage Bash
\lstdefinelanguage{dockerbash}{
  language=bash,
  keywords={sudo, docker, docker-compose, curl, cp, mkdir, rm, chmod, chown, echo, export, source, grep, sed, awk, cat, ls, cd, pwd, exit, if, then, else, fi, for, do, done, while, case, in, esac, function, return, local},
  morecomment=[l]{\#},
}

% Definition du langage Docker Compose
\lstdefinelanguage{dockercompose}{
  keywords={services, volumes, networks, image, container_name, ports, environment, depends_on, volumes, networks, restart, build, context, dockerfile, command, entrypoint, healthcheck, test, interval, timeout, retries, start_period, pull_policy, logging, driver, options, labels, dns, dns_search, extra_hosts, cap_add, cap_drop, privileged, user, working_dir, expose, links, external, name, driver_opts, ipam, config, subnet, gateway},
  keywordstyle=\color{primaryblue}\bfseries,
  sensitive=false,
  morecomment=[l]{\#},
  morestring=[b]',
  morestring=[b]",
}

% ==============================================================================
% 6. FIGURES ET TABLEAUX
% ==============================================================================
\usepackage{graphicx}                     % Inclusion d'images
\usepackage{float}                        % Placement precise des flottants
\usepackage{caption}                      % Personnalisation des legendes
\usepackage{subcaption}                   % Sous-figures et sous-tableaux
\usepackage{booktabs}                     % Tableaux professionnels
\usepackage{tabularx}                     % Tableaux adaptatifs
\usepackage{longtable}                    % Tableaux sur plusieurs pages
\usepackage{lscape}                       % Pages en mode paysage

% ==============================================================================
% 7. MATHEMATIQUES
% ==============================================================================
\usepackage{amsmath}                      % Environnements mathematiques
\usepackage{amssymb}                      % Symboles mathematiques

% ==============================================================================
% 8. SCHEMAS ET DIAGRAMMES (TikZ)
% ==============================================================================
\usepackage{tikz}
\usetikzlibrary{shapes}
\usetikzlibrary{arrows}
\usetikzlibrary{positioning}
\usetikzlibrary{calc}
\usetikzlibrary{fit}
\usetikzlibrary{backgrounds}

% ==============================================================================
% 9. ENCADRES ET BOITES
% ==============================================================================
\usepackage{fancybox}                     % Boites avec bordures

% -- Environnement infobox : encadre bleu pour informations ----------------
\newenvironment{infobox}{%
  \begin{center}
  \begin{Sbox}
  \begin{minipage}{0.9\textwidth}
  \color{infoblue}
  \textbf{Info :}\\
}{%
  \end{minipage}
  \end{Sbox}
  \setlength{\fboxsep}{10pt}
  \colorbox{infoblue!10}{\TheSbox}
  \end{center}
}

% -- Environnement warningbox : encadre jaune pour avertissements ----------
\newenvironment{warningbox}{%
  \begin{center}
  \begin{Sbox}
  \begin{minipage}{0.9\textwidth}
  \color{warnyellow!80!black}
  \textbf{Attention :}\\
}{%
  \end{minipage}
  \end{Sbox}
  \setlength{\fboxsep}{10pt}
  \colorbox{warnyellow!15}{\TheSbox}
  \end{center}
}

% ==============================================================================
% 10. EN-TETES ET PIEDS DE PAGE
% ==============================================================================
\usepackage{fancyhdr}

% Style fancy pour les pages de chapitre
\fancypagestyle{fancy}{%
  \fancyhf{}
  \fancyhead[L]{\small\textcolor{gray}{LockBits -- Dossier d'Architecture}}
  \fancyhead[R]{\small\textcolor{gray}{Projet Annuel 2025-2026}}
  \fancyfoot[C]{\thepage}
  \renewcommand{\headrulewidth}{0.4pt}
  \renewcommand{\footrulewidth}{0pt}
}

% Style plain pour les pages de debut de chapitre
\fancypagestyle{plain}{%
  \fancyhf{}
  \fancyfoot[C]{\thepage}
  \renewcommand{\headrulewidth}{0pt}
  \renewcommand{\footrulewidth}{0pt}
}

% Application du style fancy par defaut
\pagestyle{fancy}

% ==============================================================================
% 11. FORMATAGE DES TITRES
% ==============================================================================

% Chapitre : huge, bold, primaryblue
\titleformat{\chapter}[display]
  {\normalfont\huge\bfseries\color{primaryblue}}
  {\chaptertitlename\ \thechapter}
  {20pt}
  {\Huge}

% Section : Large, bold, primaryblue
\titleformat{\section}
  {\normalfont\Large\bfseries\color{primaryblue}}
  {\thesection}
  {1em}
  {}

% Sous-section : large, bold, primaryblue!80!black
\titleformat{\subsection}
  {\normalfont\large\bfseries\color{primaryblue!80!black}}
  {\thesubsection}
  {1em}
  {}

% ==============================================================================
% 12. LISTES
% ==============================================================================
\usepackage{enumitem}
\setlist{nosep, leftmargin=2em}           % Listes compactes

% ==============================================================================
% 13. HYPERLIENS ET METADONNEES PDF
% ==============================================================================
\usepackage[hidelinks]{hyperref}

\hypersetup{
  pdftitle={Dossier d'Architecture -- LockBits},
  pdfauthor={Clery A-Ferradou, Swann Kechar, Soltane Afellah, Quentin Labourdette},
  pdfsubject={Architecture SI -- Projet Annuel ESGI 2025-2026},
  pdfkeywords={EDR, Docker, FastAPI, PHP, Monitoring, CI/CD, Sécurité},
  pdfcreator={LaTeX},
  pdfproducer={pdflatex},
  colorlinks=false,
}

% ==============================================================================
% FIN DU PREAMBULE
% ==============================================================================
\endinput
